\documentclass[11pt]{article}
\usepackage[a4paper,margin=1in]{geometry}
\usepackage{booktabs}
\usepackage{longtable}
\usepackage{array}
\usepackage{hyperref}
\title{R Workshop: Curriculum}
\author{Applied R for Statistical Teaching, Research and Reproducible Analysis}
\date{}
\begin{document}
\maketitle

\section*{Positioning}
This workshop provides three hours of guided instruction per day, supported by additional workbook sections for independent study. It briefly revisits statistical theory where necessary to connect equations, assumptions and interpretation to implementation in R.

\section*{Delivery model}
Each workbook section should be labelled as Core workshop activity, Optional demonstration, or Reference/take-home material. The live sessions should prioritise the continuous capstone built around one anchor dataset.

\section*{Capstone progression}
\begin{longtable}{p{0.16\linewidth}p{0.74\linewidth}}
\toprule
Day & Capstone progress \\
\midrule
Day 2 & Select question, clean data and document variables \\
Day 3 & Produce exploratory figures and initial observations \\
Day 4 & Perform one inferential analysis \\
Day 5 & Fit and diagnose the principal model \\
Day 6 & Implement or examine part of the model mathematically \\
Day 7 & Integrate, render, review and present \\
\bottomrule
\end{longtable}

\section*{Anchor dataset principle}
Use one anchor dataset from Day 2 through Day 7: clean it, explore it, infer from it, model it, examine part of the model mathematically, and publish the complete analysis reproducibly. Label any other dataset as a demonstration dataset.

\section*{Assessment evidence}
Use the Workshop Evidence Log to record daily artefacts, peer checks and capstone progress. Repeat a short baseline task from Day 1 on Day 7 to provide evidence of participant development.

\end{document}
